%\documentclass[reprint,amsmath,amssymb,aps,twoside]{revtex4-2} % for editor use only

% For initial submission use these lines
\documentclass[amsmath,amssymb,aps,endfloats,linenumbers]{revtex4-2}



% DO NOT CHANGE THINGS BELOW HERE
\usepackage{graphicx}
\usepackage{amsmath,amssymb,amsfonts}
\usepackage{dcolumn}
\usepackage{bm}
\usepackage{siunitx}
%\usepackage{tikz,pgfplots}
\sisetup{separate-uncertainty=true, multi-part-units=single}
\usepackage[colorlinks,allcolors=blue]{hyperref}
\usepackage{cleveref}
\crefname{equation}{}{}
\crefname{figure}{Fig.}{Figs.}
\crefname{table}{Table}{Tables}
\usepackage{svg}
\svgpath{{./figures}}




% set PDF metadata
% AUTHORS EDIT THIS PART
\hypersetup{%
pdftitle={Energy is Conserved in Spring-like Poppers}, % full title here
pdfauthor={Kelly Su, Julia Bawar, Julia Khabinskiy, and Sejal Nagrani}, % full author list here
}


\usepackage{fancyhdr}
\pagestyle{fancy}
\fancyhf{}
\fancyhead[RE,RO]{J S\&E \textbf{2}, xx--xx (2026)} % FOR EDITOR ONLY
\fancyhead[LO]{Lennon \emph{et al}} % if many authors
%\fancyhead[LO]{John and Dee} % if short enough to list all authors
\fancyhead[LE]{She loves you, yeah, yeah, yeah} % full or short title depending on what fits
%\fancyhead[LE}{Several species of small furry animals gathered together in a cave and grooving with a pict} % might be too long
%\fancyhead[LE]{Several species grooving} % might be a short version of the title that fits
\fancyfoot[C]{\thepage}
\fancypagestyle{mytitlepage}{
\fancyhf{}
\fancyhead[C]{Journal of Science \& Engineering \textbf{2}, xx--xx (2026)} % FOR EDITOR ONLY
\fancyfoot[C]{\thepage}
}


\begin{document}
%\setcounter{page}{67} % FOR EDITOR USE ONLY
%\linenumbers % FOR EDITOR USE ONLY

% AUTHORS EDIT THIS PART
\title{Energy is Conserved in Spring-like Poppers}
\author{Kelly Su}
\altaffiliation{Manalapan High School} % use for non S&E authors
\author{Julia Bawar}
\author{Julia Khabinskiy}
\author{Sejal Nagrani}
\email{Contact author: 427ksu@frhsd.com} % use FRHSD address
% For S&E authors
\affiliation{Science \& Engineering Magnet Program, \href{https://manalapan.frhsd.com/}{Manalapan High School}, Englishtown, NJ 07726 USA}
\date{\today}


% AUTHORS PUT YOUR MANUSCRIPT IN BELOW HERE
\begin{abstract}
This experiment investigates energy conservation in spring-like popper toys. A popper was compressed and released to obtain its maximum height, its velocity right after hitting the maximum height, and its velocity upon release for five trials. The kinetic energy equation and the potential energy equation were evaluated with the obtained values, where the average of the yielded kinetic energy and potential energy values, 0.06894 J and 0.07922 J, respectively, were compared and found to be approximately equal. The elastic potential energy stored in the popper before launch was 0.257 J, further supporting the idea that energy is conserved. A paired t-test comparing the average kinetic energy at launch and the average gravitational potential energy at maximum height for a popper for five trials yielded a p-value of p = 0.424 > α = 0.05, revealing no statistically significant difference between the kinetic energy at launch and the potential energy at the popper’s highest point. This means that the kinetic energy at launch and the potential energy at maximum height are close enough to be comparable, further supporting the idea that energy is conserved. 
\end{abstract}

% For revtex format this will not print on the article page but will be used in setting 
% the journal metadata so include keywords. Lowercase unless they are acronyms or proper nouns
% and pick keywords that will actually help people find your work properly. 
\keywords{kinematics, pop, time series}

\maketitle\thispagestyle{mytitlepage}

%The current column width is \the\columnwidth

\section{Introduction}
Kinetic energy (KE) is the energy inside a moving object that makes it move, and can be calculated by the kinetic energy equation, which relates mass m and velocity v to produce the amount of kinetic energy K in a system[1]:

\begin{equation}
\sum \vec{F} = m \vec{a}. %K = 1/2mv^2 please :) and it says "Eq. (1)" next to it
\label{eq:fma} % example label for use with cleveref
\end{equation}

Gravitational potential energy (GPE) is the energy stored inside an object at rest, and is based on the object's vertical position away from the ground. It can be calculated by the gravitational potential energy equation, which relates the initial gravitational potential energy U0 , mass m, gravitational constant g = 9.8 m/s2, and vertical position y to produce the amount of potential energy U in a system[1]:

\begin{equation}
\sum \vec{F} = m \vec{a}. %U = U0 + mgy please :) and it says "Eq. (2)" next to it
\label{eq:fma} % example label for use with cleveref
\end{equation}

It can be hypothesized that energy is conserved in a frictionless environment, meaning that the total energy of the system should remain constant. Friction was neglected to simplify the experiment, as it has no significantly large impact on the data. The kinetic energy of the system just as the popper left the ground, and the gravitational potential energy of the system at the popper's maximum height were calculated with Eq. (1) and Eq. (2), respectively, where the system consists of the popper.

The null hypothesis and alternative hypothesis are listed below:

Null Hypothesis: KE = GPE
Alternative Hypothesis: KE ≠ GPE

The null hypothesis states that energy is conserved during the popper’s motion, meaning the kinetic energy at launch is equal to the gravitational potential energy at maximum height. This would imply that there is no significant difference between the kinetic energy at launch and the potential energy at the popper’s maximum height.

The alternative hypothesis states that energy is not conserved during the popper’s motion, meaning the kinetic energy at launch is not equal to the gravitational potential energy at maximum height.

An app called FizziQ was used to compute the popper’s velocity as it left the initial position y = 0 for Eq. (1) based on video data[2]. The initial gravitational potential energy U0 was defined as 0 for all calculations to represent the absence of energy in the system at rest. Since the conservation of energy indicates that there can only be as much kinetic energy as there was potential energy, it was deduced that, if these two values are equal when friction is negligible, energy is conserved.





\section{Methods and materials}

%figure 1 please -kelly
\begin{figure}
\begin{center}
%Image file would go here 
% \includegraphics{figures/fig1.png}
\includegraphics{figures/newfig1.pdf}
% \includesvg{figures/fig2.svg}
\end{center}
\caption{This is a caption. Captions for figures go below the figure.}
\label{fig:foo}
\end{figure}

The popper (Liberty Imports “Jumping Gens” poppers, Florham Park, NJ) is a yellow, spring-like toy that consists of a foam head with a conically structured netting attached to it, shown in FIGURE 1, with a height of 0.085 m and a weight of 0.0055 kg. The measuring device consisted of two meter sticks that were taped together so that the bottom stick would end at the 0.5 m mark of the top stick. This assembly was propped up by a group member's hands, with the edge of the ruler staying flat on the ground. Then, a popper was pinched between another group member's thumb and index finger at the side of the foam head to compress it toward y = 0, the popper's snap-through point. To release the popper, the thumb and index finger were simultaneously removed from the foam head, allowing the popper to spring upwards and reach its maximum height before moving downwards. For a total of five trials, the same person released each popper with this method for each trial to avoid confounding variables. The popper's maximum height from each trial was documented through visual observation and later, more precisely observed in videos of each trial taken at 60 frames per second with an iPhone 15. These videos were then processed through FizziQ, where each frame was handpicked, and the popper's position was indicated with a point and positioned by a group member[2]. The experiment was conducted inside a classroom in a secluded area with the windows closed, unaffected by external factors and elements, such as other individuals and wind. 

%figure 2 please -kelly
\begin{figure}
\begin{center}
%Image file would go here 
% \includegraphics{figures/fig1.png}
\includegraphics{figures/newfig1.pdf}
% \includesvg{figures/fig2.svg}
\end{center}
\caption{This is a caption. Captions for figures go below the figure.}
\label{fig:foo}
\end{figure}

To measure the exact force in the popper before release, a scale was placed under the popper to measure the weight in grams, and set to 0 to avoid counting the popper's weight, as shown in FIGURE 2. The scissor jack was placed next to the scale so that the platform sat on the popper's foam head. The scissor jack was used to adjust to a specific compression height for each trial, allowing for more exact calculations of force measurements. Additionally, the use of this tool avoided confounding variables like inconsistent human force or human error when handling the popper. With this setup, the scissor jack's platform was lowered, thereby compressing the popper to find more precise measurements of force without human error, shown in FIGURE 3. The weight was later converted to Newtons for easier calculations.



        
\section{Results}

%figure 3 please -kelly
\begin{figure}
\begin{center}
%Image file would go here 
% \includegraphics{figures/fig1.png}
\includegraphics{figures/newfig1.pdf}
% \includesvg{figures/fig2.svg}
\end{center}
\caption{This is a caption. Captions for figures go below the figure.}
\label{fig:foo}
\end{figure}

FIGURE 3, created with the Google Sheets software represents the spring force generated by different compression distances. From this, the elastic potential energy stored in the popper before release was calculated as the area under the curve. At the snap-through point, a distance of 0.05 m compressed, the area under the curve was approximately 0.257 J, according to FIGURE 3.

%figure 4 please -kelly
\begin{table}
\begin{center}
\begin{ruledtabular}
\begin{tabular}{cc}
$x$ & $\bar{x}$ \\
\colrule
0 & 1 \\ 
1 & 0 \\
\end{tabular}
\end{ruledtabular}
\end{center}
\caption{This is a table. Captions for tables go above the table.}
\label{tab:bar}
\end{table}
    
    
    
    
    
\section{Discussion}
The discussion is where you interpret the meaning of your results and whether they support your hypotheses or refute alternative hypotheses. You should refer to tables, figures, etc by number and use them in supporting your arguments here. To refer to a figure do like this: \cref{fig:foo} shows the third Zonklar moment as a function of $t$. To refer to a table like \cref{tab:bar}, do like this. 

    
    
    
    
    
    
    
\section{Acknowledgements}
Keep it short or the editors will shorten it for you. JL wrote and PM wrote the lyrics. GH contributed rhythm guitar tracks. RS contributed drumbeats. Do not write long essays about Ms Komitas, Dr E, or the S\&E program here; we don't need or want fluff. 



% References – Do not spend too much time trying to perfectly format per MLA or whatever because journals often have their own preferred formats; as long as you have author, title, journal, year, volume, page number it will be findable. Peer reviewed journals or published work only; Food Girl’s Blog is not a peer reviewed journal; Wikipedia is not a peer reviewed journal; Webster’s Dictionary is not necessary in a scientific journal article. If I find fake references your paper will be rejected. 






\bibliography{example.bib}
%References
%[1] P.A. Tipler and G. Mosca. Physics for Scientists and Engineers, 5th ed. (W H Freeman and Company, New York, 2004).
%
%“De Caelo : Aristotle : Free Download, Borrow, and Streaming : Internet Archive.” Internet Archive, 2026, archive.org/details/decaelo00aris. Accessed 12 Jan. 2026.
%
%Galilei, G. Dialogues Concerning Two New Sciences (1638).
%
%Chazot, Christophe. FizziQ. 5.0.4, Trapèze. Digital, 2025, Apple App Store.


\end{document}
